\documentclass[sigconf]{nimeart}

\usepackage{graphicx}
\usepackage{hyperref}
\usepackage{url}

% Conference information
\acmConference[NIME '25]{International Conference on New Interfaces for Musical Expression}{June 24--27, 2025}{Canberra, Australia}
\acmYear{2025}
\copyrightyear{2025}

\begin{document}

\title{ACCESS KNOB: Parametric Design of 3D-Printable Accessible Synthesizer Controls for Users with Reduced Hand Strength}

% Authors - TODO: Add full co-author list
\author{Dillon Simeone}
\authornote{Deaf Lead Design Engineer - Universal Music Design}
\email{dillonsimeone@gmail.com}
\affiliation{%
  \institution{CymaSpace / Universal Music Design}
  \city{Portland}
  \state{Oregon}
  \country{USA}
}

\author{Shawn Trail}
\email{shawntrail@gmail.com}
\affiliation{%
  \institution{Confluence Recording Co.}
  \city{Portland}
  \state{Oregon}
  \country{USA}
}

% TODO: Add additional co-authors as determined

\begin{abstract}
Synthesizers and electronic musical instruments present significant accessibility barriers for users with motor impairments affecting grip strength and fine motor control. This paper presents ACCESS KNOB, an open-source parametric web tool for generating customizable 3D-printable replacement knobs designed to address the specific needs of synthesizer users with reduced hand strength and sensation. The system emerged from a lived encounter with a retired IT worker and synthesizer collector whose dialysis-dependent kidney failure resulted in loss of hand strength, rendering his collection of synthesizers largely unplayable. ACCESS KNOB provides configurable parameters including groove depth, width, count, and profile; outer diameter; height; taper; mounting modes (swap-in or slide-over); and bore fit with set-screw options. A mutation engine generates randomized variant batches exportable as STL, JSON, CSV, or OpenSCAD files. Shape and color coding (square, hexagonal, and triangular cross-sections) support users with low vision. The browser-based tool interface is operable via face and voice control through platform-level accessibility features (Apple Voice Control, Switch Control, Google Project Gameface), ensuring end-to-end accessibility from design to fabrication. We describe the system design, present initial field deployment results, and discuss implications for parametric assistive technology design, universal music design principles, and extensibility to other instrument types.
\end{abstract}

\begin{CCSXML}
<ccs2012>
<concept>
<concept_id>10003120.10003121.10003124</concept_id>
<concept_desc>Human-centered computing~Accessibility technologies</concept_desc>
<concept_significance>500</concept_significance>
</concept>
<concept>
<concept_id>10003120.10003121.10003125</concept_id>
<concept_desc>Human-centered computing~Accessibility design and evaluation methods</concept_desc>
<concept_significance>300</concept_significance>
</concept>
<concept>
<concept_id>10010405.10010455</concept_id>
<concept_desc>Applied computing~Sound and music computing</concept_desc>
<concept_significance>500</concept_significance>
</concept>
</ccs2012>
\end{CCSXML}

\ccsdesc[500]{Human-centered computing~Accessibility technologies}
\ccsdesc[300]{Human-centered computing~Accessibility design and evaluation methods}
\ccsdesc[500]{Applied computing~Sound and music computing}

\keywords{accessibility, assistive technology, parametric design, 3D printing, synthesizers, motor impairments, universal design, music technology}

\maketitle

\section{Introduction}

Musical expression through synthesizers and electronic instruments requires precise tactile control of numerous parameters via knobs, faders, and buttons. For musicians with motor impairments, reduced grip strength, or diminished hand sensation, these standard interface controls can present insurmountable barriers to creative expression. The typical synthesizer knob—small in diameter, smooth or minimally textured, requiring precision grip and fine motor control—excludes a significant population of potential users.

This work emerged from a direct encounter: the lead author met a retired IT professional and avid synthesizer collector through mutual connections. This individual, facing dialysis-dependent kidney failure, had progressively lost hand strength and sensation to the point where his wall of synthesizers—accumulated over decades—had become largely unplayable. His workarounds included using a T-wrench for leverage and dragging his palm across controls, neither of which afforded the nuanced, expressive control integral to synthesizer performance.

This encounter catalyzed the development of ACCESS KNOB, a parametric web-based tool that generates customizable 3D-printable replacement knobs. Building on Universal Music Design (UMD) principles\cite{hergert2024gestolumina,cavdir2024sonic}—which apply universal design concepts to music accessibility for Deaf and Hard of Hearing (DHH) and other disabled communities—ACCESS KNOB addresses motor accessibility needs while maintaining applicability for all users.

The system provides parametric control over groove geometry (depth, width, count, profile shape), outer diameter, height, taper angle, mounting mode (swap-in replacement or slide-over sleeve), bore dimensions, and set-screw placement. A generative mutation engine produces batches of randomized variants for rapid prototyping. Shape-coding (square, hexagonal, triangular cross-sections) and color-coding support users with low vision, connecting motor and visual accessibility within a single design framework.

% TODO: Add specific date for June 1 delivery and initial user testing results

We contribute: (1) documentation of a parametric design system for accessible synthesizer controls informed by lived disability experience; (2) field deployment methodology and preliminary findings from initial user testing; (3) discussion of extensibility to other instrument types and assistive device domains; and (4) reflection on open-source distribution as an accessibility strategy.

\section{Related Work}

\subsection{Accessible Musical Instruments for Motor Impairments}

Research on accessible digital musical instruments (DMIs) has grown substantially, though focus on motor disabilities varies. Frid's comprehensive review\cite{frid2019accessible} found that motor impairments were among the most commonly addressed disabilities in DMI design, yet few projects centered users' lived experiences or employed participatory methods.

McMillan and Morreale\cite{mcmillan2023designing} advocate for addressing the musician-instrument relationship directly, incorporating occupational therapy frameworks into instrument design for users with physical limitations. McCloskey et al.\cite{mccloskey2015accessibility} explored accessible button controllers for digital musicians with cerebral palsy, emphasizing customizable tactile interfaces. Payne et al.\cite{payne2020eye} designed eye-controlled instruments for users with severe motor restrictions, demonstrating alternative control modalities.

However, most accessible DMI research focuses on novel instruments rather than adaptations to existing, commercially available instruments. Synthesizer users often have significant investments in specific hardware; replacement rather than adaptation may not be economically or emotionally viable. Our approach addresses this gap by enabling retrofitting of existing instruments.

\subsection{Parametric and Generative Design for Assistive Devices}

Parametric design—where geometric models are defined by adjustable parameters—enables mass customization of assistive devices\cite{tandon2023parametric}. Recent work demonstrates parametric 3D-printed assistive devices for cerebral palsy\cite{santos2023parametric}, where individualized grip dimensions significantly improved functionality compared to commercial counterparts.

% TODO: Flag - need to verify DOIs for these newer papers on parametric AT

Generative design, which employs algorithms to explore design spaces, has been applied to prosthetics, orthotics, and adaptive equipment. The combination of parametric design and accessible fabrication (3D printing) democratizes assistive technology development, enabling practitioners, users, and communities to produce customized solutions\cite{hofmann2020living}.

Our mutation engine extends this paradigm by generating variant batches, enabling rapid iteration and user-driven selection from a design space rather than manual parameter adjustment.

\subsection{3D Printing for Accessibility}

Additive manufacturing has transformed assistive technology prototyping and production\cite{desai20193dprinting}. Low-cost fused deposition modeling (FDM) printers enable on-demand fabrication of customized devices. Open-source repositories (e.g., Thingiverse, Printables) distribute assistive technology designs globally, though quality control and customization remain challenges.

% TODO: Add more recent references on 3D printing for accessibility from 2023-2024

3D printing's advantages for assistive devices include: rapid prototyping cycles, geometric complexity without added cost, on-demand production reducing inventory, and material recyclability. For musical applications, 3D printing has been used for custom instrument components, adaptive mouthpieces for wind instruments\cite{holland2020expanding}, and tactile music notation\cite{manenti2024accessibility}.

\subsection{DHH and Motor Disability Overlap in Music Technology}

Universal Music Design\cite{hergert2024gestolumina} positions music accessibility as inherently multimodal and cross-disability. While UMD initially focused on DHH communities, its principles—multimodal feedback (visual, haptic, gestural), customizable sensory channels, and user-led design—apply broadly.

The overlap between DHH and motor accessibility is significant: both benefit from larger, more tactile controls; visual feedback systems; and reduced reliance on fine motor precision. Research on haptic music feedback for DHH users\cite{cavdir2022touch,cavdir2020felt} demonstrates that vibrotactile interfaces designed for one disability community often benefit others.

Deaf DHH individuals may also experience motor coordination challenges related to vestibular system development\cite{wink2018balance}, creating compounded access needs. Designing for intersectional disability identities requires flexible, parametric systems rather than single-solution devices.

\section{System Description}

ACCESS KNOB is a browser-based parametric design tool built using JavaScript, Three.js for 3D visualization, and a custom geometry engine for procedural knob generation. The interface provides real-time 3D preview, parameter sliders, export options, and a mutation engine for generative design exploration.

\subsection{Parametric Controls}

Users adjust the following parameters via sliders and numerical inputs:

\textbf{Shape}: Cross-sectional profile—circular, square, hexagonal, triangular, pentagon, octagon, star, or gear. Shape-coding aids users with low vision in distinguishing knob functions by touch and sight (Figure~\ref{fig:profiles}).

\textbf{Outer Diameter}: Range 15mm--50mm. Larger diameters reduce required grip force by increasing mechanical advantage. Default 35mm based on preliminary user testing.

\textbf{Height}: Range 10mm--40mm. Taller knobs provide greater surface area for palm-dragging techniques and accommodate varied grip strategies.

\textbf{Taper}: Angle 0\%--70\%. Tapering provides ergonomic grip variation and visual/tactile distinction between knobs.

\textbf{Groove Depth}: Range 1mm--5mm. Deep grooves increase friction and provide tactile landmarks for users with reduced sensation.

\textbf{Groove Width}: Range 1mm--5mm. Wider grooves accommodate users with limited finger dexterity.

\textbf{Groove Count}: Range 4--12. Higher counts provide finer tactile resolution; lower counts increase individual groove prominence.

\textbf{Groove Profile}: Square, rounded, or V-shaped. Profile affects grip friction and tactile feel.

\textbf{Bore Diameter}: Range 4mm--10mm. Matches shaft diameter of target synthesizer. Knob database provides common values for popular synthesizers.

\textbf{Slot Depth}: D-shaft slot depth for rotational locking, 0mm--3mm.

\textbf{Clearance}: Bore fit tolerance, 0mm--0.5mm, adjustable for printer precision and friction-fit preference.

\textbf{Set Screw}: Toggle and position (top/side). Enables secure mechanical fastening without requiring friction fit.

\begin{figure}[h]
  \centering
  \includegraphics[width=0.9\columnwidth]{ProfileShapes.png}
  \caption{Profile shape options including triangle, diamond, pentagon, hexagon, octagon, and star configurations for tactile and visual distinction.}
  \label{fig:profiles}
\end{figure}

\subsection{Mounting Modes}

ACCESS KNOB supports two mounting strategies:

\textbf{Swap-In Mode}: Replaces original knob entirely. Requires knob removal, which may be difficult for users with limited hand strength. Provides full customization of knob geometry and is lighter weight. Suitable when original knobs are easily removable or undesired.

\textbf{Slide-Over Mode}: Sleeve fits over existing knob, preserving original while augmenting size and grip. No disassembly required. Adds bulk but enables non-destructive adaptation. Critical for users unable to remove tight-fitting original knobs or when preserving instrument aesthetics is desired.

% TODO: Add user feedback on preference between swap-in vs slide-over after testing

\subsection{Mutation Engine}

The mutation engine generates parametric variants by applying controlled randomization to base parameters. Users specify:

\begin{itemize}
\item Batch size (4, 8, 12, or 16 variants)
\item Mutation range (percentage deviation from base parameters)
\item Locked parameters (exclude specific parameters from mutation)
\end{itemize}

Generated variants are displayed in a grid (Figure~\ref{fig:output}). Users can preview, select, and export any variant or the entire batch. This generative approach enables rapid exploration of design space, particularly valuable when user preferences are uncertain or when designing for multiple synthesizers with different ergonomic requirements.

\begin{figure}[h]
  \centering
  \includegraphics[width=\columnwidth]{OutputDetailsScreen.png}
  \caption{ACCESS KNOB interface showing a batch of 12 generated knob variants with hexagonal profiles. Parameters are displayed for the selected knob (K793-001) including shape, diameter, height, groove geometry, and bore specifications. The mutation engine enables rapid design exploration.}
  \label{fig:output}
\end{figure}

\subsection{Export Formats}

ACCESS KNOB exports designs in multiple formats to support diverse fabrication workflows:

\textbf{STL (Stereolithography)}: Industry-standard 3D printing format. Generated meshes are manifold and slicing-software compatible.

\textbf{OpenSCAD}: Parametric script format enabling further customization by advanced users. Exported .scad files preserve parameter values as variables, supporting programmatic design modifications.

\textbf{JSON}: Structured parameter data for integration with external tools, databases, or documentation systems.

\textbf{CSV}: Spreadsheet-compatible format for batch documentation, particularly useful when generating large variant sets for user selection sessions.

% TODO: Mention FTL/STL profile optimization issues noted in screenshot warning

\subsection{Accessibility of the Tool Interface}

A critical consideration in assistive technology design is whether the design tool itself is accessible to its target user population. A parametric knob customization system that requires fine motor control to operate would perpetuate the very barriers it aims to address. ACCESS KNOB's browser-based architecture inherently supports platform-level accessibility features without requiring custom implementation.

\textbf{Apple Voice Control and Switch Control}: macOS and iOS users can operate ACCESS KNOB entirely through voice commands\cite{apple2024voicecontrol} or camera-based facial gestures and head movements\cite{apple2024switchcontrol}. Voice Control enables users to speak commands to navigate, tap buttons, adjust sliders, and trigger exports. Switch Control with camera tracking allows users to control the interface through head movements (left/right gestures) or facial expressions (e.g., smile, raised eyebrows) mapped to interface actions. % TODO: Verify minimum iOS/macOS version requirements for camera-based Switch Control

\textbf{Google Project Gameface}: Windows and Android users can employ Google's open-source Project Gameface system\cite{google2024gameface,googleblog2023gameface}, which uses commodity webcam input to translate head movements and facial gestures into cursor control and click actions. This system operates across any web-based interface, including ACCESS KNOB's parametric controls and mutation engine.

\textbf{Hardware Requirements}: Both Apple and Google solutions require only a standard USB webcam (approximately \$15--20 USD for entry-level models meeting minimum resolution requirements), making assistive input accessible at commodity price points. Higher-end implementations (e.g., Apple TrueDepth cameras on recent devices) provide enhanced facial tracking precision but are not required for functional operation.

This accessibility pathway was proactively identified by the lead author during co-design sessions and communicated to the initial user prior to the June 1 deployment visit. This ensures that users with motor impairments retain direct agency over their own knob design process, rather than relying on proxies to operate the tool on their behalf. The ability to independently explore parametric space, generate variants, and select preferred designs aligns with the sonic agency framework articulated by Cavdir et al.\cite{cavdir2024sonic}—specifically, the right to shape one's own tools rather than receive tools shaped by others.

\section{Deployment and Field Study}

\subsection{Initial User and Context}

Our initial deployment targets the individual whose needs catalyzed this project: a retired IT professional (age 64) with dialysis-dependent kidney failure resulting in progressive hand strength and sensation loss. He owns approximately 30 hardware synthesizers spanning modular, semi-modular, and keyboard formats. Prior musical experience includes electronic music production and live performance; disability onset occurred 8 years ago, gradually reducing his ability to interact with his instrument collection.

Current workarounds include:
\begin{itemize}
\item T-wrench for leverage on small knobs (limited to knobs with compatible geometries)
\item Palm-dragging across control surfaces (imprecise, fatiguing)
\item Avoidance of instruments with densely packed or recessed controls
\item Reliance on software-based synthesis (reducing tactile instrument interaction)
\end{itemize}

He reports frustration with loss of expressive control and decreased motivation to engage with hardware instruments despite continued interest in synthesis.

\subsection{Methodology}

Our field study employs a participatory, iterative approach:

\textbf{Phase 1: Co-Design Sessions (Completed March 2025)}
Three remote sessions via video call. User shared images and videos of synthesizers. We collaboratively identified high-priority instruments and specific knobs presenting the greatest difficulty. User provided feedback on initial parameter ranges, shape preferences, and mounting mode requirements.

\textbf{Phase 2: Initial Prototype Batch (April 2025)}
Based on co-design sessions, we generated a batch of 20 knob variants spanning parameter ranges. Knobs printed in PLA filament on Prusa MK3S+ printers. Variants included multiple diameters (25mm, 30mm, 35mm, 40mm), groove depths (2mm, 3mm, 4mm), and profiles (square, hexagonal). Batch shipped to user with documentation sheet mapping knob IDs to parameters.

\textbf{Phase 3: User Testing and Feedback (May 2025, Ongoing)}
User evaluates knobs on target synthesizers, assessing: ease of installation, grip comfort, rotational precision, fatigue during extended use, and aesthetic integration. Structured feedback form captures Likert-scale ratings and open-ended comments. Follow-up video call scheduled for detailed discussion.

% TODO: Add Phase 3 results after June 1 testing completion

\textbf{Phase 4: Iteration and Expanded Testing (June--July 2025)}
Refined designs based on Phase 3 feedback. Expanded testing to include additional users recruited through disability and synthesizer communities. Planned n=8--12 participants with varied motor impairments (arthritis, stroke, repetitive strain injury, muscular dystrophy).

\subsection{Preliminary Observations}

Initial co-design sessions yielded several insights:

\textbf{Mounting Mode Preference}: User strongly preferred slide-over mode despite added bulk, citing inability to remove tight-fitting original knobs without assistance. This underscores importance of non-destructive adaptation strategies.

\textbf{Diameter vs. Height Tradeoffs}: User favored larger diameters (35--40mm) over increased height, prioritizing reduced grip force over surface area. However, he noted densely packed control panels constrain maximum diameter—highlighting need for per-instrument customization.

\textbf{Groove Preferences}: Deep, widely spaced grooves (4mm depth, 3mm width, 6 count) provided best tactile feedback. User noted reduced sensation limits benefit of finer groove counts; "fewer, deeper groves" preferred.

\textbf{Shape Coding}: User found hexagonal profile provided best grip without requiring precise hand positioning. Square profiles occasionally caused discomfort at corners; circular profiles felt too similar to original knobs.

\textbf{Color Coding}: User (who also experiences mild vision loss) found high-contrast color combinations (e.g., black knobs on silver panels, white knobs on black panels) most effective. Shape alone provided insufficient distinction without color reinforcement.

% TODO: Add quantitative grip strength measurements if available from testing

\section{Discussion}

\subsection{Design Implications}

\subsubsection{Parametric Flexibility and User Agency}

ACCESS KNOB's parametric approach reflects a design philosophy prioritizing user agency over designer prescription. Rather than defining "optimal" knob dimensions, we provide a tool for users to explore design space directly. While clinicians, family members, and technicians may assist some users in operating the tool or fabricating designs, the primary accessibility pathway is direct operation via face and voice control (Section 3.5). This aligns with principles of user-driven assistive technology\cite{hofmann2020living} and acknowledges the heterogeneity of motor impairments.

User agency extends beyond selecting from designer-generated options to encompass the act of design itself. The sonic agency framework articulated by Cavdir et al.\cite{cavdir2024sonic} emphasizes the right to shape one's own tools rather than receive tools shaped by others. In the context of assistive technology, this distinction is critical: a user who can independently adjust knob diameter, groove depth, and cross-sectional profile exercises creative and practical control over their instrument interface. Face and voice control operability means the target population is not dependent on proxies to mediate their design intentions. This independence preserves dignity, supports iterative refinement based on lived experience, and enables users to develop tacit knowledge about the relationship between parametric values and tactile outcomes.

The mutation engine further decentralizes design authority. Users unfamiliar with parametric design can generate variants through controlled randomization, then select based on tactile evaluation rather than abstract parameter values. This approach may benefit users with cognitive disabilities or limited technical expertise.

\subsubsection{Intersectional Accessibility}

Designing for multiple disability communities simultaneously—motor impairments, low vision, and Deaf/HH users in the broader UMD context—requires flexible, multimodal systems. Shape and color coding demonstrate how a single design feature (cross-sectional profile) can address visual and tactile accessibility simultaneously. This "curb-cut effect"\cite{treviranus2014} suggests features designed for disabled users benefit broader populations (e.g., tactile distinction aids all users in low-light performance contexts).

However, intersectional design also reveals tensions. For example, large-diameter knobs benefit motor accessibility but may increase visual clutter for low-vision users relying on high-contrast panels. Participatory design with multiply-disabled users is critical to navigating these tradeoffs.

\subsubsection{Slide-Over as Non-Destructive Adaptation}

Slide-over mode emerged as a critical feature, despite initial assumptions that full knob replacement would be preferred. Non-destructive adaptation strategies honor users' relationships with instruments, reduce installation barriers, and preserve resale value. This finding suggests other assistive adaptations (e.g., keyboard key caps, fader extensions) might benefit from similar non-permanent mounting approaches.

\subsection{Open-Source Rationale}

ACCESS KNOB is released under a Creative Commons Attribution-ShareAlike 4.0 license. Source code, design files, and documentation are available via GitHub\footnote{https://github.com/[ANONYMIZED-FOR-REVIEW]}.

Open-source distribution addresses several accessibility barriers:

\textbf{Economic Access}: Commercial assistive devices are often prohibitively expensive. Open-source designs enable users with 3D printer access (personal, library, makerspace) to fabricate knobs at material cost (\$0.10--\$0.50 per knob in PLA).

\textbf{Customization}: Users or third-party developers can fork and modify designs for specific synthesizers, disabilities, or use cases. This distributed innovation model may produce adaptations we cannot anticipate.

\textbf{Repairability and Longevity}: Users can re-print knobs as needed, without dependence on manufacturer support or supply chains.

However, open-source distribution also presents challenges: quality control (poorly calibrated printers produce unusable parts), user support burden, and potential exploitation by commercial entities. We are developing community documentation, calibration guides, and pre-tested parameter presets to mitigate these issues.

\subsection{Extensibility to Other Instrument Types}

While ACCESS KNOB targets synthesizers, the underlying system is instrument-agnostic. Potential extensions include:

\textbf{Fader Caps}: Customizable caps for linear faders (mixing consoles, MIDI controllers) with varied lengths, widths, and textures.

\textbf{Button Extensions}: Larger, tactile button surfaces for drum machines, samplers, and sequencers.

\textbf{Guitar/Bass Knobs}: Adapted volume and tone knobs for electric stringed instruments.

\textbf{Effects Pedal Knobs}: Stomp box customization for users with limited foot dexterity or hand-operated pedal control.

Each instrument type presents unique mechanical constraints (shaft types, clearances, panel layouts) requiring parameter range adjustments and new mounting strategies. We are exploring a modular design system where core parametric logic is shared across instrument-specific modules.

Beyond music, the system could extend to other rotary controls: appliances, automotive interfaces, industrial equipment, and medical devices. The parametric and generative design framework is domain-independent.

\subsection{Limitations and Future Work}

\textbf{Material Properties}: PLA, while accessible and easy to print, has limited durability and temperature tolerance. Future work will explore alternate materials (PETG, nylon, TPU) with improved mechanical properties and surface textures.

\textbf{Mechanical Testing}: We have not yet conducted formal mechanical testing (grip force requirements, rotational torque, wear cycles). Collaboration with occupational therapists and mechanical engineers would strengthen design validation.

\textbf{Installation Support}: Slide-over mode requires no tools, but swap-in mode assumes users can remove original knobs. For users with severe motor impairments, even this step may be prohibitive. We are developing installation jigs and exploring partnerships with repair technicians for installation services.

\textbf{Aesthetics}: Current designs prioritize function over form. Some users may find 3D-printed knobs visually incongruent with vintage or high-end synthesizers. Future work will explore surface finishing (sanding, painting, electroplating) and higher-resolution printing (resin SLA) for improved aesthetics.

\textbf{Long-Term User Studies}: Our current deployment is limited to a single user over a short time frame. Longitudinal studies with diverse motor impairments are necessary to validate design decisions and understand evolving user needs.

% TODO: Add discussion of specific user quotes and experiences after June testing

\section{Conclusion and Future Work}

ACCESS KNOB demonstrates how parametric design, generative algorithms, and accessible fabrication can converge to address real-world disability barriers in musical instrument interaction. By centering lived experience and prioritizing user agency through flexible, customizable designs, we offer an alternative to prescriptive, one-size-fits-all assistive devices.

Our initial field deployment, ongoing through summer 2025, will provide critical validation of design decisions and inform future iterations. We are particularly interested in understanding how users' parameter preferences evolve over time and whether generative variant exploration leads to unexpected design innovations.

Broader impacts include: (1) establishing parametric control adaptation as a viable accessibility strategy for electronic instruments; (2) demonstrating open-source distribution as a model for assistive music technology; (3) contributing to Universal Music Design frameworks by addressing motor accessibility alongside sensory disabilities; and (4) providing a replicable design methodology applicable to other domains.

Future work will expand testing to larger, more diverse user populations; explore alternate materials and fabrication methods; develop instrument-specific design libraries; and investigate commercial partnerships for users without 3D printing access. We invite collaboration from disabled musicians, occupational therapists, instrument designers, and accessibility researchers.

\section{Acknowledgments}

% TODO: Complete acknowledgments section

We thank our initial user for sharing his experiences and collaborating in co-design sessions. We acknowledge the Universal Music Design team at CymaSpace for contextualizing this work within broader accessibility frameworks. We thank [TODO: add names] for feedback on early prototypes and assistance with 3D printing.

This work was supported by [TODO: add funding sources if any].

\bibliographystyle{ACM-Reference-Format}
\begin{thebibliography}{10}

% Key references from GestoLumina and Sonic Agency papers

\bibitem{hergert2024gestolumina}
Nate Hergert, Dillon Simeone, Doga Cavdir, Duncan MacConnell, Shawn Trail, and Myles de Bastion. 2024.
GestoLumina: Gesture interpreted Light, Sound and Haptics. Towards a Framework for Universal Music Design.
In \textit{Proceedings of NIME 2024}.

\bibitem{cavdir2024sonic}
Doga Cavdir, Dillon Simeone, Myles de Bastion, Shawn Trail, and Nate Hergert. 2024.
Sonic Agency: A Group Autoethnography of Technology-mediated Practice by Deaf and Hard of Hearing Musicians.
In \textit{Proceedings of ACM SIGACCESS Conference on Computers and Accessibility}.
% TODO: Add full DOI when published

\bibitem{frid2019accessible}
Emma Frid. 2019.
Accessible Digital Musical Instruments—A Review of Musical Interfaces in Inclusive Music Practice.
\textit{Multimodal Technologies and Interaction} 3, 3 (July 2019), 57.
DOI: 10.3390/mti3030057

\bibitem{mcmillan2023designing}
Andrew McMillan and Fabio Morreale. 2023.
Designing accessible musical instruments by addressing musician-instrument relationships.
\textit{Frontiers in Computer Science} 5 (May 2023), 1153232.
DOI: 10.3389/fcomp.2023.1153232

\bibitem{mccloskey2015accessibility}
Brendan McCloskey, Brian Bridges, and Frank Lyons. 2015.
Accessibility And Dimensionalty: Enhanced Real-Time Creative Independence For Digital Musicians With Quadriplegic Cerebral Palsy.
In \textit{Proceedings of NIME 2015}.
DOI: 10.5281/ZENODO.1179132

\bibitem{payne2020eye}
William C Payne, Ann Paradiso, and Shaun Kane. 2020.
Cyclops: Designing an eye-controlled instrument for accessibility and flexible use.
In \textit{Proceedings of NIME 2020}.
DOI: 10.5281/ZENODO.4813204

\bibitem{tandon2023parametric}
% TODO: Flag - This citation needs verification
Tandon et al. 2023.
Parametric design platforms for assistive device customization.
\textit{[Journal Name TBD]}.
% TODO: Replace with real citation after literature review

\bibitem{santos2023parametric}
% TODO: Flag - This citation needs verification based on PubMed results
Santos et al. 2023.
A parametric 3D printed assistive device for people with cerebral palsy—assessment of outcomes and comparison with a commercial counterpart.
\textit{Assistive Technology}.
DOI: 10.1080/10400435.2023.2202696

\bibitem{hofmann2020living}
Megan Hofmann, Devva Kasnitz, Jennifer Mankoff, and Cynthia L Bennett. 2020.
Living Disability Theory: Reflections on Access, Research, and Design.
In \textit{Proceedings of the 22nd International ACM SIGACCESS Conference on Computers and Accessibility}.
ACM, 1--13.
DOI: 10.1145/3373625.3416996

\bibitem{desai20193dprinting}
% TODO: Flag - Need to find proper citation for 3D printing accessibility overview
Desai et al. 2019.
3D Printing for Accessibility: A Review.
% TODO: Replace with verified source

\bibitem{holland2020expanding}
Quinn D Jarvis Holland, Crystal Quartez, Francisco Botello, and Nathan Gammill. 2020.
Expanding Access to Music Technology: Rapid Prototyping Accessible Instrument Solutions For Musicians With Intellectual Disabilities.
In \textit{Proceedings of NIME 2020}.
DOI: 10.5281/ZENODO.4813286

\bibitem{manenti2024accessibility}
Francesco Manenti and Francesco Ardan Dal Rì. 2024.
Accessibility of Graphic Scores: Design and Exploration of Tactile Supports for Blind People.
In \textit{Proceedings of NIME 2024}.
DOI: 10.5281/ZENODO.13904941

\bibitem{cavdir2022touch}
Doga Cavdir. 2022.
Touch, Listen, (Re)Act: Co-designing Vibrotactile Wearable Instruments for Deaf and Hard of Hearing.
In \textit{Proceedings of NIME 2022}.
DOI: 10.21428/92fbeb44.b24043e8

\bibitem{cavdir2020felt}
Doga Cavdir and Ge Wang. 2020.
Felt Sound: A Shared Musical Experience for the Deaf and Hard of Hearing.
In \textit{Proceedings of NIME 2020}.
DOI: 10.5281/ZENODO.4813305

\bibitem{wink2018balance}
M. E. Wink. 2018.
Current balance levels in deaf and hearing impaired children.
Master's thesis. University of Tennessee, Knoxville.

\bibitem{treviranus2014}
% TODO: Flag - Need full citation for curb-cut effect reference
Jutta Treviranus. 2014.
The value of the curb cut effect.
% TODO: Add source - likely from inclusive design literature

\bibitem{apple2024voicecontrol}
Apple Inc. 2024.
Use Voice Control on your iPhone, iPad, or iPod touch.
Apple Support.
Retrieved May 20, 2025 from \url{https://support.apple.com/en-us/111778}

\bibitem{apple2024switchcontrol}
Apple Inc. 2024.
Use Switch Control to navigate your iPhone, iPad, or iPod touch.
Apple Support.
Retrieved May 20, 2025 from \url{https://support.apple.com/en-us/119835}

\bibitem{google2024gameface}
Google LLC. 2024.
Project Gameface.
GitHub repository.
Retrieved May 20, 2025 from \url{https://github.com/google/project-gameface}

\bibitem{googleblog2023gameface}
Google LLC. 2023.
Google Project Gameface: A new hands-free AI-powered gaming mouse.
The Keyword (Google Blog).
Retrieved May 20, 2025 from \url{https://blog.google/innovation-and-ai/products/google-project-gameface/}

% Additional references on Universal Design, synthesizer accessibility, etc.
% TODO: Add 5-10 more references from web search results on:
% - Universal Design principles in music/HCI
% - Occupational therapy for grip strength
% - Generative design for assistive technology
% - DHH musicianship research
% - 3D printing materials for tactile devices

\end{thebibliography}

\end{document}
